\documentclass[12pt]{article}
\usepackage[T1]{fontenc}
\usepackage[utf8]{inputenc}
\usepackage{lmodern}
\usepackage{textcomp}
\usepackage{enumitem}
\usepackage{geometry}
\geometry{margin=1in}
\begin{document}
\begin{center}
Answer Key - Economics - Undergraduate Level Assessment 4 \
\small (Answers are ordered exactly as in the question paper.)
\end{center}

\section*{Multiple Choice}
\begin{enumerate}[label=\arabic*.]
\item \textbf{Answer:} A
\\ \textit{Options:} A) fallen 50 percent.; B) fallen 25 percent.; C) risen 15 percent.; D) risen 20 percent.; E) fallen 20 percent.
\\ \textit{Note:} The correct answer is incorrect because an increase in the Consumer Price Index (CPI) from 120 to 150 indicates a rise in prices by 25 percent, not a fall, as the CPI measures inflation by comparing the current price level to a base level.
\item \textbf{Answer:} A
\\ \textit{Options:} A) keep the price constant as demand is elastic.; B) decrease price as demand is inelastic.; C) decrease the quality of the product.; D) increase price as demand is elastic.; E) increase price as demand is inelastic.
\\ \textit{Note:} The correct answer is based on the concept of price elasticity of demand, where a value greater than 1 indicates elastic demand, meaning that keeping the price constant will lead to an increase in quantity demanded and subsequently raise total revenue.
\item \textbf{Answer:} E
\\ \textit{Options:} A) Scarcity is relevant only for non-renewable resources in the circular flow model; B) Scarcity is only relevant in the factor markets of the circular flow model; C) Scarcity is only considered in the household sector of the circular flow model; D) Scarcity is addressed through government intervention in the circular flow model; E) Scarcity is implicit in the entire analysis of the circular flow model
\\ \textit{Note:} The circular flow model illustrates the interactions between households and firms in an economy, inherently incorporating scarcity by demonstrating how limited resources are allocated to meet unlimited wants through the flow of goods, services, and money.
\item \textbf{Answer:} D
\\ \textit{Options:} A) increase long-run output with no significant impact on inflation.; B) have no significant impact on either output or inflation.; C) lower the unemployment rate while also lowering the rate of inflation.; D) increase short-run output but it is the source of long-run inflation.; E) increase the nation's long-run capacity to produce.
\\ \textit{Note:} The quantity theory of money posits that while an increase in the money supply can temporarily boost short-run output by stimulating demand, it ultimately leads to higher prices and inflation in the long run as the economy adjusts to the increased monetary base without a corresponding increase in real output.
\item \textbf{Answer:} A
\\ \textit{Options:} A) Historical data contradicts theory; B) Oligopolies are less common in practice than in theory; C) There is no contradiction between theory and practice; D) Historical data is insufficient to determine the efficiency of oligopolies; E) Oligopolies typically demonstrate higher innovation rates than predicted
\\ \textit{Note:} The correct answer is supported by the fact that traditional economic theory often assumes that oligopolies operate efficiently by maximizing profits and minimizing costs, while historical data shows that many oligopolistic industries have engaged in practices such as price-fixing and reduced competition, leading to inefficiencies that contradict the theoretical expectations.
\end{enumerate}

\section*{True / False}
\begin{enumerate}[label=\arabic*.]
\setcounter{enumi}{5}
\item \textbf{Answer:} False
\\ \textit{Options:} A) True; B) False
\\ \textit{Note:} The statement is false because real GDP is calculated by adjusting nominal GDP for inflation using the GDP deflator, which means real GDP would be lower than \$6000 when the deflator is 200.
\item \textbf{Answer:} True
\\ \textit{Options:} A) True; B) False
\\ \textit{Note:} Raising taxes on disposable income reduces the amount of money consumers have available to spend, thereby decreasing overall consumer demand and helping to mitigate the effects of built-in inflation.
\item \textbf{Answer:} False
\\ \textit{Options:} A) True; B) False
\\ \textit{Note:} Social overhead capital refers to public investments in infrastructure and essential services, such as transportation and utilities, rather than private sector investments in technology startups.
\item \textbf{Answer:} False
\\ \textit{Options:} A) True; B) False
\\ \textit{Note:} Lower interest rates typically encourage private investment by making borrowing cheaper, which can complement expansionary fiscal policy rather than enhance its impact negatively.
\item \textbf{Answer:} True
\\ \textit{Options:} A) True; B) False
\\ \textit{Note:} In the long run, a tax cut in a full employment economy will lead to increased consumption and investment, which can enhance productivity and efficiency, ultimately resulting in higher real output without affecting the price level since the economy is already at full capacity.
\end{enumerate}

\section*{Long Form Response}
\begin{enumerate}[label=\arabic*.]
\setcounter{enumi}{10}
\item \textbf{Answer:} The relationship between income elasticity of demand for real money and real income growth can be analyzed using the given elasticity of 0.8 and a demand growth rate of 2.6\%. The income elasticity of demand indicates how responsive the quantity demanded of real money is to changes in real income. To find the growth rate of real income, we can use the formula that links these variables: growth rate of real income = demand growth rate / income elasticity. Substituting the values, we calculate it as 2.6\% / 0.8, which results in a growth rate of approximately 3.47\% per year for real income. This illustrates that as the demand for real money grows, it is essential to understand the proportionate increase in real income necessary to support that demand.
\\ \textit{Note:} correct because it accurately applies the formula linking income elasticity of demand to real income growth by dividing the demand growth rate by the income elasticity, demonstrating how changes in demand for real money are proportional to changes in real income.
\item \textbf{Answer:} Public debt can create inflationary pressures in an economy primarily by enhancing consumption spending. When a government borrows money, it often spends this capital on various programs and services, which can stimulate demand in the economy. This increase in consumption can lead to higher prices, especially if the economy is operating near its capacity. Additionally, public debt can represent a backlog in purchasing power, as individuals and businesses may anticipate future tax increases or inflation that could erode their savings. As a result, delayed purchasing decisions may occur, leading to a surge in demand when consumers finally spend their accumulated resources, further exacerbating inflationary trends. Overall, the dynamics of public debt can significantly influence consumption patterns and overall price levels in an economy.
\\ \textit{Note:} correct because it illustrates how government borrowing can boost consumption spending, thereby increasing demand and potentially leading to inflation, while also acknowledging the future implications for purchasing power due to anticipated tax increases.
\end{enumerate}

\vfill
\noindent\textit{End of Answer Key}
\end{document}